% ============================================================================
%  APARTADO: Pipeline Computacional — de la Simulación UrQMD al VAE
%  Sugerencia de ubicación: Capítulo 4 "Construcción del Análisis",
%  como nuevas subsecciones dentro de la Sección "Estado Actual: Simulaciones
%  MCG + UrQMD" (después de la línea que hoy termina en \label{eq:elbo} /
%  antes de \chapter{Resultados}), o alternativamente como cuerpo del
%  Capítulo 6 "Avances y Trabajo a Realizar" en sustitución de la lista
%  genérica de "Etapa I/II/III", ya que lo que sigue describe el trabajo
%  YA implementado, no solo lo planeado.
% ============================================================================

\subsection{Generación de los Eventos UrQMD}
\label{sec:generacion_urqmd}

\noindent
La línea de base del análisis se construye a partir de $10^{5}$ eventos de
colisiones Bi+Bi a $\sqrt{s_{NN}} = 11$ GeV generados con UrQMD en modo
mínimamente sesgado (\textit{minimum bias}), es decir, sin restringir a priori
el parámetro de impacto $b$, de modo que el muestreo cubre todo el espectro
de centralidades. Esta energía se eligió porque corresponde a la región del
diagrama de fases donde, de acuerdo con las predicciones discutidas en la
Sección~\ref{sec:CEP}, se localiza el punto crítico (CEP), maximizando la
densidad bariónica neta accesible en NICA.

\noindent
Cada uno de los $10^{5}$ eventos se generó con una semilla aleatoria
(\textit{seed}) distinta. Esto es una condición necesaria, no solo una
cuestión de reproducibilidad: puesto que el observable central del análisis
es la \emph{fluctuación evento por evento} de $\langle p_T \rangle$
(Ecuación~\ref{eq:momento_n_muestral}), una correlación espuria entre
eventos generados con semillas repetidas o encadenadas subestimaría
artificialmente la varianza $\sigma^2_{\langle p_T \rangle}$ y, con ella,
cualquier señal dinámica atribuible al CEP. Por razones de capacidad de
cómputo, la corrida completa se fraccionó en varios lotes
(\texttt{Bi\_11GeV\_MB1}, \dots, \texttt{Bi\_11GeV\_MB5}), cada uno con su
propio conjunto de subdirectorios \texttt{run\_NNNN/}; todos los scripts de
análisis descritos más adelante aceptan múltiples carpetas como entrada y
las combinan en una sola pasada de lectura, de forma que la fragmentación es
transparente para el análisis. Los archivos de salida, en formato
\texttt{.f14} (estándar de UrQMD), se alojaron en almacenamiento en la nube
para permitir su descarga y procesamiento desde distintos equipos de
cómputo.

% ------------------------------------------------------------------
\subsection{Lectura de los Archivos \texttt{.f14}: \texttt{read\_f14.py}}
\label{sec:read_f14}

\noindent
Se desarrolló un lector propio, \texttt{read\_f14.py}, para el formato ASCII
de salida de UrQMD 3.4 (19 columnas por partícula: posición y tiempo de
\textit{freeze-out}, cuadrimomento $(E,p_x,p_y,p_z)$, masa, código de
partícula \texttt{ityp}, isospín, carga \texttt{chg}, y, de manera crucial
para este trabajo, el contador de colisiones \texttt{ncl}). A partir del
cuadrimomento se derivan, partícula por partícula y directamente al nivel
del generador (es decir, \emph{antes} de cualquier simulación del detector),
las cuatro cantidades cinemáticas usadas en todo el análisis posterior:

\begin{equation}
    p_T=\sqrt{p_x^2+p_y^2},\quad
    P=\sqrt{p_x^2+p_y^2+p_z^2},\quad
    \eta=\tfrac{1}{2}\ln\!\frac{P+p_z}{P-p_z},\quad
    \phi=\operatorname{atan2}(p_y,p_x).
\end{equation}

\noindent
El parámetro de impacto $b$ de cada evento se extrae de la línea de
cabecera (\texttt{"i"}) de UrQMD. La variable \texttt{ncl} permite
clasificar cada partícula según su historial de colisiones, clasificación
que se usa de manera sistemática en todos los scripts posteriores:

\begin{itemize}
    \item $\texttt{ncl}=0$: \textbf{espectadoras} — nucleones que no sufrieron
    ninguna colisión y por tanto conservan (aproximadamente) su cuadrimomento
    original.
    \item $\texttt{ncl}=1$: \textbf{primarias} — participaron en exactamente
    una colisión.
    \item $\texttt{ncl}>1$: \textbf{secundarias} — productos de colisiones o
    reacciones múltiples (resonancias, rescattering).
    \item $\texttt{ncl}>0$ (primarias $\cup$ secundarias): agrupadas bajo la
    etiqueta \textbf{participantes} en los análisis de fluctuaciones.
\end{itemize}

\noindent
Tanto la lectura de un archivo (\texttt{iter\_events}) como la de un
directorio completo de corridas (\texttt{iter\_all}) están implementadas
como generadores de Python: cada evento se procesa y se libera de memoria
antes de leer el siguiente, de modo que el consumo de RAM es constante
($\mathcal{O}(1)$) sin importar si se procesan cien o cien mil eventos. Esta
decisión de diseño es la que permite que todos los scripts de la sección
siguiente escalen al conjunto completo de $10^{5}$ eventos en un único
recorrido (\textit{streaming}) de los datos.

% ------------------------------------------------------------------
\subsection{Extracción de Observables desde \texttt{.f14}}
\label{sec:pipeline_observables}

\noindent
Sobre la base de \texttt{read\_f14.py} se construyó una batería de scripts
de análisis, cada uno enfocado en un observable físico distinto:

\begin{itemize}

\item \textbf{\texttt{build\_dataset.py}} — convierte el flujo de eventos en
un archivo HDF5 compacto, con histogramas por evento de $p_T$ (60 bins,
0--3 GeV/$c$), $\eta$ (80 bins, $-8$ a $8$) y $\phi$ (72 bins, $-\pi$ a
$\pi$), más el vector de parámetros de impacto $b$. La escritura se realiza
por lotes (\textit{batches}) con compresión \texttt{gzip} y \textit{datasets}
redimensionables, preservando el mismo consumo constante de memoria de
\texttt{read\_f14.py}. Este archivo HDF5 es la representación de entrada
que consume directamente la VAE de histograma completo
(\texttt{vae\_pt.py}, Sección~\ref{sec:vae}).

\item \textbf{\texttt{plot\_observables.py}} (versión final; la primera
iteración, \texttt{p\_o.py}, procesaba una sola carpeta y una sola especie a
la vez) — genera los espectros $\frac{1}{N_{\text{ev}}}\frac{dN}{dp_T}$,
$\frac{1}{N_{\text{ev}}}\frac{dN}{d\eta}$ y $\frac{1}{N_{\text{ev}}}\frac{dN}{d\phi}$
combinando varias carpetas de simulación en una sola pasada. Para cada
observable se generan tres gráficas: \emph{Todas} (sin desglose de especie),
\emph{Cargadas} ($\pi^{+},\pi^{-},K^{+},K^{-},p$, identificados por
\texttt{ityp} y el signo de \texttt{chg}) y \emph{No cargadas}
($\pi^{0},K^{0},\bar{K}^{0},n$ — el neutrón solo aparece en este último
grupo, ya que es la única especie neutra masiva de interés), y dentro de
cada gráfica se superponen las curvas de partículas \emph{participantes}
y \emph{espectadoras}. El script admite además la opción
\texttt{--by-species} para obtener una gráfica individual por cada una de
las 11 especies. Estas son las figuras presentadas en la
Sección~\ref{sec:pt_eta_resultados} (Figuras~\ref{fig:pt_pions}
a~\ref{fig:pt_charged}).

\item \textbf{\texttt{fluc\_observables.py}} — clasifica cada evento en una
de 10 clases de centralidad según cortes en $b$ (siguiendo el mapeo
Glauber $b\!\to\!$centralidad usado en el resto del trabajo, de 0--10\%
hasta 90--100\%) y produce, por clase: la distribución del número de
partículas por evento, un diagrama de caja tipo \textit{vela japonesa}
(mínimo/Q1/mediana/Q3/máximo/media) de la multiplicidad, la distribución
de $p_T$ normalizada por evento, y la distribución de $\langle p_T\rangle$
por evento. Es el primer diagnóstico visual de que la dispersión de
$\langle p_T\rangle$ efectivamente varía con la centralidad, antes de
calcular los momentos formales.

\item \textbf{\texttt{pt\_moments.py}} — calcula, para cada clase de
centralidad, los momentos muestrales $\mu_1=\langle\langle p_T\rangle\rangle$
(Ecuación~\ref{eq:grand_mean_pt}), $\mu_2=\sigma^2_{\langle p_T\rangle}$ y
$\mu_3$ (Ecuaciones~\ref{eq:momento_n_muestral}
y~\ref{eq:skewness_muestral}) de la distribución $P(\langle p_T\rangle)$, y
los contrasta con la predicción cerrada de la distribución Gamma derivada
en la Sección~\ref{sec:momentos_pt_gamma}: $\alpha=\mu_1^2/\mu_2$,
$\beta=\mu_2/\mu_1$, $\mu_{3,\text{pred}}=\mu_2^2/\mu_1$ y
$S_{\text{pred}}=2/\sqrt{\alpha}$. Este script es, en la práctica, la
validación numérica de la derivación analítica del Capítulo 5: compara el
\textit{skewness} medido contra el predicho por el modelo Gamma en función
de la centralidad. El error de $\mu_1$ se reporta como error estándar de la
media; los de $\mu_2$, $\mu_3$ y $S$, por \textit{bootstrap} (500
remuestreos), al no poder asumirse normalidad para esos estimadores.

\item \textbf{\texttt{pt\_correlator.py}} y \textbf{\texttt{pt\_correlator\_norm.py}}
— calculan el correlador de dos partículas $\langle \Delta p_{T,i}\,\Delta
p_{T,j}\rangle$ (Ecuación~\ref{eq:correlador_2p}, equivalente a
$\sigma^2_{\text{dyn}}$ de la Ecuación~\ref{eq:sigma2_dyn}) y su versión
normalizada $\sqrt{\langle N\rangle\,\langle \Delta p_{T,i}\,\Delta
p_{T,j}\rangle}/\langle\langle p_T\rangle\rangle$
(cf. Ecuación~\ref{eq:cv_correlador}), separando las partículas por
\texttt{ncl} en primarias y secundarias en vez de por cortes de
$|\eta|<1.0$, que se usaron en una versión previa y se retiraron
deliberadamente por decisión metodológica. El error se estima por
\textit{jackknife} \textit{delete-1} en lugar de \textit{bootstrap}, pues
la razón $\sqrt{\langle N\rangle\cdot\text{corr}}/\mu_1$ amplifica el
ruido cuando el correlador es pequeño, y el \textit{jackknife} resulta más
estable en ese régimen.

\end{itemize}

\noindent
En conjunto, estos cinco scripts comparten la misma capa de lectura
(\texttt{read\_f14.py}), los mismos cortes de centralidad y las mismas
convenciones de selección de partículas, de modo que los observables que
alimentan el análisis estadístico (Capítulo 4) y los que alimentan la VAE
(siguiente subsección) provienen de exactamente los mismos filtros
aplicados a los mismos $10^{5}$ eventos.

% ------------------------------------------------------------------
\subsection{Nota sobre la Interpretación del Espectro de $\eta$: Espectadoras
y Aceptancia del MPD}
\label{sec:nota_eta_espectadoras}

\noindent
Un resultado que conviene precisar es el comportamiento de las partículas
espectadoras en la variable $\eta$: su distribución se concentra en una
región relativamente plana para $|\eta|\lesssim 4$, en contraste con las
partículas participantes, que dominan a rapidez media. Es importante
señalar que $\eta$ se calcula, como se mostró arriba, directamente del
cuadrimomento reportado por UrQMD al nivel de generador; no depende en
absoluto de la geometría del MPD ni del tiempo en el que una partícula
alcanzaría físicamente algún subdetector, de modo que el origen de esta
estructura es puramente cinemático y no instrumental. La explicación física
es la siguiente: para $\sqrt{s_{NN}}=11$ GeV la rapidez del haz es
$y_{\text{haz}}=\operatorname{arccosh}\!\big(\sqrt{s_{NN}}/2m_N\big)\approx
2.45$ (Ecuación~\ref{eq:rapidity}); los nucleones espectadores, al no
colisionar, conservan una rapidez cercana a $y_{\text{haz}}$, con un
ensanchamiento adicional producido por el movimiento de Fermi intrínseco de
los nucleones dentro del núcleo antes del choque. Esa combinación es la que
produce la meseta observada hasta $|\eta|\approx 4$, en contraste con la
región de rapidez media ($|y|<1.2$) donde se concentran las partículas
participantes y que constituye la aceptancia central del MPD
(Sección~\ref{cap:NICA}).

\noindent
La relevancia de esta observación para el diseño experimental es, entonces,
de \emph{aceptancia} y no de dinámica de colisión con el detector: dado que
la TPC del MPD cubre únicamente $|\eta|<1.5$, las partículas espectadoras
caen sistemáticamente fuera de su aceptancia y deben medirse con un
subdetector dedicado a ángulos muy pequeños respecto al eje del haz —el ZDC
(\textit{Zero-Degree Calorimeter})—, que es precisamente el instrumento que
el MPD utiliza para calibrar la centralidad y el plano de reacción a partir
de la energía de las espectadoras. En otras palabras: la separación
espacial entre participantes y espectadoras que se observa en estos
histogramas de $\eta$ es la razón física por la que el MPD necesita
subdetectores distintos (TPC/TOF centrales frente a ZDC frontal) para cada
población, y no una consecuencia de la forma cilíndrica del detector
actuando sobre trayectorias ya generadas.

% ==================================================================
\subsection{Implementación del Autoencoder Variacional (VAE)}
\label{sec:vae}

\noindent
Sobre la base teórica de la Sección~\ref{sec:vae_teoria} (función objetivo
ELBO, Ecuación~\ref{eq:elbo}) se implementaron y compararon cuatro
representaciones de entrada distintas para el VAE, cada una encapsulada en
su propio script. La comparación no es arbitraria: cada representación
responde a una crítica metodológica distinta sobre qué información del
evento debe ver la red, y las últimas versiones incorporan directamente la
retroalimentación recibida durante el desarrollo.

\begin{enumerate}

\item \textbf{\texttt{vae\_pt.py} — histograma completo de $p_T$.} La
entrada es el histograma normalizado de $p_T$ del evento (60 bins,
0--3 GeV/$c$, el mismo \textit{binning} que produce
\texttt{build\_dataset.py}), es decir, la VAE ve la \emph{forma completa}
del espectro. Arquitectura: codificador $60\to256\to128\to(\mu_z,\log
\sigma_z^2)$, decodificador $z_{\dim}\to128\to256\to60$ con salida
\texttt{softmax} (reconstruye una distribución de probabilidad). Función de
pérdida $\mathcal{L}=\text{MSE}(x,\hat{x})+\beta D_{\text{KL}}$ (o
entropía cruzada binaria con la bandera \texttt{--bce}). El espacio latente
se evalúa correlacionándolo (Pearson) con $b$, la clase de centralidad,
$\langle p_T\rangle$ y la multiplicidad, y mediante \textit{clustering}
no supervisado (K-means y DBSCAN, evaluados con \textit{silhouette score}).

\item \textbf{\texttt{vae\_pt\_mean.py} — histograma \textit{bootstrap} de
$\langle p_T\rangle$.} En lugar de usar el histograma crudo de $p_T$ (que
mezcla la forma del espectro con ruido estadístico de conteo finito), este
script construye, para cada evento $i$, un histograma \emph{derivado}: se
remuestrean con reemplazo las $N_i$ partículas del evento $K$ veces, se
calcula $\langle p_T\rangle$ en cada remuestreo, y el histograma normalizado
de esas $K$ estimaciones es el vector de entrada. El ancho de ese
histograma tiene una lectura física directa: es la incertidumbre
estadística de $\langle p_T\rangle_i$ dado un $N_i$ finito
($\sim\sigma_{p_T}/\sqrt{N_i}$), de modo que la VAE ve simultáneamente la
posición \emph{y} la precisión de la estimación. La multiplicidad mínima
$N_{\min}$ ya no se fija a mano: se estima mediante un estudio de
\textit{bootstrap} independiente sobre el error estándar de
$\langle p_T\rangle$ en función de $N$ (función \texttt{estimate\_n\_min}).
Todo el análisis —carga, cálculo de $N_{\min}$, entrenamiento y
\textit{clustering}— se realiza por separado para cada una de las 10 clases
de centralidad, entrenando un modelo independiente por clase.

\item \textbf{\texttt{vae\_moments.py} (y su variante \texttt{vae\_3m.py})
— tres momentos intra-evento.} Reduce cada evento a un vector de solo tres
escalares, $x=(\mu_1,\mu_2,\mu_3)$, calculados sobre las partículas
\emph{dentro} de ese evento:
$\mu_1=\langle p_T\rangle_i$, $\mu_2=\frac{1}{N}\sum(p_{T,j}-\mu_1)^2$,
$\mu_3=\frac{1}{N}\sum(p_{T,j}-\mu_1)^3$. Conviene subrayar, para que quede
explícito en la tesis, que estos $\mu_2,\mu_3$ describen la forma del
espectro \emph{dentro} de un solo evento y no deben confundirse con
$\sigma^2_{\langle p_T\rangle}$ ni con el \textit{skewness} $S$ del
\emph{ensemble} $P(\langle p_T\rangle)$ definidos en las
Ecuaciones~\ref{eq:momento_n_muestral} y~\ref{eq:skewness_muestral}: son
observables físicos válidos por derecho propio, pero conceptualmente
distintos. Al tratarse de una entrada de solo 3 dimensiones, la arquitectura
se reduce drásticamente ($3\to64\to32\to(\mu_z,\log\sigma_z^2)$).

\item \textbf{\texttt{vae\_moments\_V3.py} — momentos como desviaciones
respecto a la centralidad, con correcciones metodológicas.} Es una
revisión de \texttt{vae\_moments.py} que incorpora, en respuesta a
retroalimentación recibida durante el desarrollo:
\begin{itemize}
    \item $d_z\le 2$: no tiene sentido comprimir una entrada de dimensión 3
    a un espacio latente de dimensión mayor o igual.
    \item Corte $N\ge N_{\min}=10$ para garantizar validez estadística de
    $\mu_2$ y $\mu_3$, con $N_{\min}$ eventualmente determinado también por
    \textit{bootstrap} (misma idea que en \texttt{vae\_pt\_mean.py}) en
    vez de fijarse arbitrariamente.
    \item Entrada redefinida como \emph{desviaciones} respecto al promedio
    de la clase de centralidad, $\delta\mu_k=\mu_k-\langle\mu_k\rangle_{\text{clase}}$,
    de forma que la red aprenda las fluctuaciones intrínsecas del evento y
    no la dependencia geométrica trivial con la centralidad.
    \item Arquitectura reducida ($3\to16\to8\to d_z$) para evitar
    sobreparametrización, y $\beta\ge0.5$ para prevenir el colapso del
    posterior.
    \item Un \texttt{StandardScaler} \emph{independiente por clase de
    centralidad}, en vez de uno global: dado que la varianza de las
    fluctuaciones $\delta\mu_k$ escala con la multiplicidad (y por tanto con
    la centralidad), un escalador único distorsiona la magnitud relativa
    de las fluctuaciones centrales frente a las periféricas
    (heterocedasticidad, diagnosticada explícitamente antes de aplicar la
    corrección).
\end{itemize}

\end{enumerate}

\noindent
En los cuatro casos, la representación latente $\mathbf{z}$ se somete al
mismo protocolo de validación: correlación de Pearson/Spearman con $b$, la
clase de centralidad y los observables de $p_T$; agrupamiento no
supervisado (K-means/DBSCAN) evaluado con \textit{silhouette score}; y
comparación entre entrada original y reconstrucción por clase de
centralidad. Esta progresión —de un histograma de 60 bins, a una
representación \textit{bootstrap} de la incertidumbre, a un resumen de solo
tres momentos, y finalmente a las desviaciones de esos momentos respecto a
su propia clase de centralidad— documenta explícitamente cómo cada
iteración responde a una limitación metodológica identificada en la
anterior, y constituye la base sobre la que se construirá la comparación
cuantitativa de la Etapa IV (Sección~\ref{cap:avances}).

\noindent
\textbf{Nota para la versión final del texto/repositorio:} el nombre y el
docstring interno de los scripts de la VAE de momentos no son
completamente consistentes entre sí —por ejemplo, \texttt{vae\_3m.py}
conserva en su docstring el nombre \texttt{vae\_moments.py}, y el
\texttt{vae\_moments.py} actual se autodescribe como una ``v4'' aun cuando
el archivo \texttt{vae\_moments\_V3.py} implementa las correcciones
etiquetadas como versión 3. Antes de referenciar estos archivos como
apéndice o material suplementario de la tesis conviene renombrarlos de
forma consecutiva (p. ej. \texttt{vae\_moments\_v1.py} \ldots
\texttt{v4.py}) y actualizar sus docstrings para que coincidan, de manera
que un lector del código (por ejemplo, el comité) pueda seguir la
cronología de las correcciones sin ambigüedad.
